\documentclass{report}

\usepackage[margin=1.5in]{geometry}

\usepackage{float}
\usepackage{graphicx}
\usepackage{hyperref}
\hypersetup{colorlinks=true,linkcolor=blue}

\usepackage{listings}
\lstMakeShortInline{|}
\lstset{columns=flexible,language=SQL,basicstyle=\ttfamily}

\def\sectionautorefname{Section}
\def\subsectionautorefname{Section}
\def\subsubsectionautorefname{Section}

\newcommand{\myref}[1]{\autoref{#1} on page \pageref*{#1}}
\newcommand{\tabref}[1]{\lstinline|#1| table (see \myref{tab:#1})}

\newcommand{\intern}{This table may be interned, see \myref{def:intern}.}

\setlength{\belowcaptionskip}{10pt}

\begin{document}

\begin{titlepage}
  \begin{center}
    \includegraphics[width=10cm]{../../lib/images/bluej-icon-256.png}\\
    \huge BlueJ Blackbox Data Collection\\~\\Researchers' Handbook
  \end{center}
\end{titlepage}

This document acts as a manual for the BlueJ data collection project,
``Blackbox''.  It is intended to explain, in detail, the schema of the
collected data, and how the data relates to the BlueJ IDE.  It assumes a basic
familiarity with BlueJ's capabilities, such as the debugger, codepad, object
bench and so on.

\begingroup
\let\clearpage\relax
\tableofcontents
\endgroup

\listoftables

\begin{table}
\caption[Possible \lstinline!master_events.name! values]{List of different
  values allowed for the \lstinline!name! column in the \lstinline!master_events! column, with
  links to sections in which that event is described.}
\label{tab:event_names}
\begin{center}
\begin{tabular}{p{6cm}@{}l}
|"assertion"| \dotfill & \myref{evt:assertion}\\
|"bench_get"| \dotfill & \myref{evt:bench_get}\\
|"bench_to_fixture"| \dotfill & \myref{evt:bench_to_fixture}\\
|"bluej_finish"| \dotfill & \myref{evt:bluej_finish}\\
|"bluej_start"| \dotfill & \myref{evt:bluej_start}\\
|"cancel_test"| \dotfill & \myref{evt:cancel_test}\\
|"codepad"| \dotfill & \myref{evt:codepad}\\
|"compile"| \dotfill & \myref{evt:compile}\\
|"debugger_breakpoint_add"| \dotfill & \myref{evt:debugger_breakpoint_add}\\
|"debugger_breakpoint_remove"| \dotfill & \myref{evt:debugger_breakpoint_remove}\\
|"debugger_close"| \dotfill & \myref{evt:debugger_close}\\
|"debugger_continue"| \dotfill & \myref{evt:debugger_continue}\\
|"debugger_halt"| \dotfill & \myref{evt:debugger_halt}\\
|"debugger_hit_breakpoint"| \dotfill & \myref{evt:debugger_hit_breakpoint}\\
|"debugger_open"| \dotfill & \myref{evt:debugger_open}\\
|"debugger_stepinto"| \dotfill & \myref{evt:debugger_stepinto}\\
|"debugger_stepover"| \dotfill & \myref{evt:debugger_stepover}\\
|"debugger_terminate"| \dotfill & \myref{evt:debugger_terminate}\\
|"edit"| \dotfill & \myref{evt:edit}\\
|"end_test"| \dotfill & \myref{evt:end_test}\\
|"file_add"| \dotfill & \myref{evt:file_add}\\
|"file_delete"| \dotfill & \myref{evt:file_delete}\\
|"fixture_to_bench"| \dotfill & \myref{evt:fixture_to_bench}\\
|"inspector_hide"| \dotfill & \myref{evt:inspector_hide}\\
|"inspector_show"| \dotfill & \myref{evt:inspector_show}\\
|"invoke_method"| \dotfill & \myref{evt:invoke_method}\\
|"package_closing"| \dotfill & \myref{evt:package_closing}\\
|"package_opening"| \dotfill & \myref{evt:package_opening}\\
|"project_closing"| \dotfill & \myref{evt:project_closing}\\
|"project_opening"| \dotfill & \myref{evt:project_opening}\\
|"remove_object"| \dotfill & \myref{evt:remove_object}\\
|"rename"| \dotfill & \myref{evt:rename}\\
|"resetting_vm"| \dotfill & \myref{evt:resetting_vm}\\
|"run_test"| \dotfill & \myref{evt:run_test}\\
|"start_test"| \dotfill & \myref{evt:start_test}\\
|"terminal_close"| \dotfill & \myref{evt:terminal_close}\\
|"terminal_open"| \dotfill & \myref{evt:terminal_open}\\
|"vcs_commit"| \dotfill & \myref{evt:vcs_commit}\\
|"vcs_history"| \dotfill & \myref{evt:vcs_history}\\
|"vcs_share"| \dotfill & \myref{evt:vcs_share}\\
|"vcs_status"| \dotfill & \myref{evt:vcs_status}\\
|"vcs_update"| \dotfill & \myref{evt:vcs_update}\\
\end{tabular}
\end{center}
\end{table}

\chapter{Concepts and Caveats}

There are several concepts involved in the Blackbox project, and it is
very important when doing analysis on the data to understand the concepts
involved, and particularly the limitations of the data collection.  If you
naively assume that one virtual user is one real person, you may draw false
conclusions from your analysis.

\section{Users}
\label{def:users}

BlueJ stores its properties file in a per-user location.  For example, on
Linux, it uses |~/.bluej/bluej.properties|, whereas on Windows 7 it is (on a typical
home installation) |C:\Users\Joe\bluej\bluej.properties|.

This properties file is the only mechanism that BlueJ has to have any sense of
identity.  If two physical people use the same account, they will appear to
BlueJ as a single user (this is still quite common on home Windows systems).
If someone uses BlueJ on a university machine, and on a laptop, and on a home
machine, they will appear to be three users.
If the computer is set up to not maintain a user's profile correctly (e.g. a
school system which wipes a user's profile on logout -- we know some systems
do this!), they will get a new profile each time they load BlueJ, and thus
will always appear to be a new user.  If the same bluej.properties file is
shared between many physical users (e.g. they all get the same profile on
login), they will all appear to be the same user (operating with many parallel sessions).

The user's unique identifier or opt-out preference is stored in this
properties file.  Thus the concept of a ``user'' is tied to the properties
file.  It is important to understand this when reasoning about users.

\section{Projects}
\label{def:projects}

BlueJ deals with projects.  A project is a directory with a BlueJ project
file, and source files (and class files, and maybe a README, and whatever else
the user stores in the directory).  In Blackbox terms, projects are associated
with particular user ids.  They are identified only by the associated user,
and their full path (hashed, for anonymity).

If a user makes a copy of a project and opens it, it will appear to be an
entirely separate project (though you could post-process the data to see that
the project has the same source code as another of the user's projects).  If a
user has the project on a USB memory stick, and it gets assigned a different
drive letter next time they plug the stick in, the project will appear to be
different.  If a person loads the project from two different machines, even if
the project has the same path on both machines, they will appear as separate
projects belonging to separate users.

\section{Packages}
\label{def:packages}

A project will have at least one package (the unnamed, default Java package),
and potentially many.  BlueJ displays a main interface window for each package
that the user opens within the project.  Thus there can be as many codepads as
there are packages in the project.  Some BlueJ features are per-package (codepad,
object bench), while others are per-project (debugger, version control).

\section{Sessions}
\label{def:sequence_id}

Each time that BlueJ is opened, we deem the time from the opening to the
closing to be ``a session''.  Each time BlueJ is loaded, the client generates
a unique session identifier to label all events from that session with.

If a user loads many projects at once in a single instance of BlueJ, these
will all count as one session.  If a user loads BlueJ twice in parallel
(e.g. by clicking the shortcut too many times), both sessions will be
recorded.  But for this reason, it is usually a good idea to use sessions and projects to
process events, rather than solely relying on event timestamps (which will
interleave events from the two sessions).  If a user accidentally closes BlueJ
and immediately re-opens it, that will count as two sessions (though you could
spot this using the timestamps on the events in post-processing).

There is no time limit on how long a session can last -- the client-generated
identifier is relied upon to differentiate sessions.  It is possible for a
session to last for weeks if a user keeps BlueJ open on their always-on
machine.  One potentially common case is that if a user never shuts down their
laptop, and instead always makes it sleep, the session can last for a long
time if they leave BlueJ open on their laptop (but see
\myref{def:interruption} for why they may not send us much data this way, if the session
stops sending data at any point).

\section{Participants}
\label{sec:participants}

All the data in the BlueJ data collection is anonymous.  Part of the intention
of the project is to support local experiments, where you, the local
researchers, collect additional information about users.  The additional
information should be held by you, and the experiment approved and conducted
according to your local ethics committee.

To allow you to link your participant data with our data, we have a mechanism
for BlueJ users to specify an experiment identifier and a participant
identifier (strings, up to 32 characters).  These two identifiers should be
assigned by you, the local researcher.  We request that you assign an
experiment identifier that is unlikely to clash with other researcher's
identifiers.  The participant identifier is entirely up to you, although we
would expect that you use a generic identifier (e.g. a number), and that you
hold the mapping between the participant's real identity and the participant
number, as is standard practice.  (That is: do not ask the participants to store their name, or
university login or similar, in the participant identifier.)

As described in section \myref{sec:views}, you will be able to retrieve
participant identifiers for a specific experiment identifier, and retrieve
session data for known experiment identifiers and participant identifiers.
This means that you will not be able to see experiment or participant
identifiers for experiments for which you do  not know the identifier.  This
is deliberate, to help protect information about experiment participation and
participant identification.

\section{Events}

The BlueJ client sends a stream of events to the BlueJ server, which
acknowledges them.  Ultimately, all the information recorded on the server
came by way of these discrete events.  Things like users or sessions are
reconstructed from the individual events that have been sent.  Every event
that the server receives, without exception, is recorded with an entry in the \tabref{master_events}.

\section{Interrupted Connections}
\label{def:interruption}

There are potential problems with missing events or entire sessions in the
event stream for a user.  You should
never assume you have a user's complete history of interaction with BlueJ.
Notably, entire sessions might be missing due to users using BlueJ without an
internet connection (e.g. on an aeroplane).  Similarly, you should not assume
that you can see all interactions with a project -- source code could be
edited outside BlueJ.

Within a session, you can assume that you have all events, sequentially
consistent, from the first \texttt{bluej\_start} event, up until the last
event that is recorded in the session.  This may not be the actual final event
that the user performed in their real session -- the user's
internet connection may drop, or BlueJ may crash or get killed before a batch of events
is sent, and so on.  But there will never be a ``hole'' in the middle of a
session where events are lost, then resumed again.  \textbf{Any event that fails to
send will cause the BlueJ client to stop sending for the rest of the session}
-- each event is only sent once the event before it in the session has succeeded.

It is a deliberate decision to not attempt to resume sending data if previous
data has failed to send.  Part of this decision is technical -- we do not want
to keep retrying a send if the user has no Internet connection (or if our own
server is inaccessible for some reason), and we do not want to cache lots of
events on the client in the hope that the connection one day re-appears.

The other reason is that analysis becomes much harder if events within a session
may be missing -- data consistency is a nightmare.  For example, you may see
that a user has closed an inspector, but not how that inspector appeared.  Or
you may see an execution take place, but not know what breakpoints are set in
the code.  If gaps are allowed in the event stream, we would have to send a
lot more state information with each event to make analysis possible.

\section{Spam}
\label{def:spam}

The BlueJ data collection project server is an open web server to which anyone
can submit, and the BlueJ client is an open-source program.  There is no
technical way to guarantee that the data coming in is from an unmodified BlueJ
client -- it could be modified, or it could simply be made-up spam data being
submitted.  It remains to be seen whether this will be an issue.

\section{Source Code Consistency}

When a package is opened, the full state of the source code of the project is
sent to the server.  Similarly, when a class is added to the project while
BlueJ is running, the full code is sent.  For all other edit events within a session, a diff is sent
between the last successfully sent source code state (during that session) and
the new state.  Using a diff minimises the amount of data sent and stored.

You can assume that the source code diffs within a session are consistent.
However, you should not assume that you have seen the full state of the source
code between sessions.  Source code may be edited outside BlueJ or may be edited
by BlueJ when the user has no Internet connection.  You should be able to
detect this by comparing the source code state between the end of one session
and the start of the next.

\section{The Debug VM}

BlueJ has two virtual machines.  One, generally known as the server VM,
handles:

\begin{itemize}
\item The graphical interface
\item Code editing
\item Compilation
\end{itemize}

No user code is ever run on the server VM.  All user code is run on a second
VM, known as the client VM or the debug VM.  (There is only ever one debug VM
at a time, for each project.)  So things like:

\begin{itemize}
\item Object construction
\item Method invocations
\item Codepad use
\item JUnit tests
\end{itemize}

Are all run on the debug VM.  When the debugger is used, the interface and
control code are in the server VM, using hooks into the debug VM to do
stepping and so on.  Whenever user code is executed, it is possible that the
result is that the debug VM terminates, but this does not mean that BlueJ has
exited.  User code that runs forever will not interrupt BlueJ, but will
prevent further use of the debug VM.


\section{Extensions}

BlueJ has an extension mechanism that allows programmatic operation of BlueJ
from third-party extensions.  (Not to mention that BlueJ is open-source,
anyway.)  The extensions can trigger many of the events described in this
document.  We do not annotate which events are triggered by extensions and
which are not.  This annotation would be a lot of extra technical work, and
probably for not much purpose, since we do not believe many BlueJ users use
extensions that will trigger these events.

We do record which extensions are loaded in the blackbox data, so you can use
this to rule out any extensions which could be problematic (or more likely, to
whitelist extensions that are not an issue).


\section{Database Structure and Notes}
\label{sec:assoc}

The database is used as a persistence layer by our Ruby on Rails server, which
uses an object-store model.  Every table in the database, except many-to-many link tables, has an
integer |id| field which acts as a primary key.  Relations in the database are
handled in two ways:

\begin{enumerate}
\item When the destination table is fixed and known -- e.g. a master event has a user --
  there is a field in the |master_events| table named |user_id| which links to
  the |users| table.
\item When the destination table may vary, e.g. a master event has a subordinate
  ``event'' record which may be from one of many tables, there is both an
  |event_id| and an |event_type| field.  The |event_type| field contains an
  uppercased camel-cased, non-underscored\footnote{This is Rails' doing -- don't ask.} version
  of the table name.  So if a particular master event links to a compile
  event, the |event_type| field will have |"CompileEvent"|, and the |event_id|
  will correspond to an |id| in the |compile_events| table.
\end{enumerate}

\subsection{Interning}
\label{def:intern}

There are some tables in the database which are likely to feature the same
information over and over again, such as the installation details.  When the
server comes to add a record to these tables with duplicate data to a previous
record, it \textit{may} reference the existing record instead, to save space, like
interning strings.  This is a performance optimisation that should have no
effect on the semantics.  Tables that do this are tagged as such.

\chapter{Master Events Table}

The master events table is the place where all events are entered.  There are
lots of different types of events that can be received, each with a variety of
different data, but all of them get one entry in the |master_events|
table, with links to other tables where necessary.

\label{tab:master_events}
\begin{table}[H]
\begin{center}
\caption[\lstinline!master_events! schema]{Schema for the \lstinline!master_events! table. 
}
\begin{tabular}{l@{\hspace{2cm}}l@{\hspace{1cm}}l}
Field Name & Type & Can Be Null?\\ \hline
\end{tabular}
\end{center}
\end{table}

The most important field in the table is the |name| field, which describes
which kind of event it is, and determines which of the other fields will be
present.  The possible values for the |name| field are given in \myref{tab:event_names}.

Each master event has links (as per \myref{sec:assoc}) to other tables:
\begin{itemize}
\item |user_id| -- present for all events
\item |session_id| -- present for all events
\item |project_id| -- present for all events except |"bluej_start"| and |"bluej_finish"|
\item |package_id| -- present for some events (see further details below).
\item |event_id|, |event_type| -- present for some events (see further details below).
\item |participant_id| -- present for all events, if the user currently has
  participant information specified.
\item |client_address_id| -- present for all events.
\end{itemize}

The |created_at| field is a server-side, trustworthy datestamp of when
the event was received.  The |source_time| is the client-side time,
which is likely (but not guaranteed) to be sequentially consistent, but should
not be trusted to be accurate.

The |sequence_num| field is provided by the client to make sure events are
received in correct sequence (within a given session).  The server ensures
that the sequence numbers are valid (i.e. that they begin at 1, and each
subsequent event increments the sequence number by 1).

\label{tab:participants}

The |participant_id| field links to a hidden |participants| table.  If you
want to investigate this information, see \myref{sec:views} on how to
access this information.

\label{tab:client_addresses}

The |client_address_id| field links to a hidden |client_address| table.
However, you can reasonably assume that there is a one-to-one mapping between
|client_address_id| and IP addresses from which the event is sent.


The subsequent sections in this chapter describe the different possible values for the
|master_events.name| field, and what that means for the other fields,
especially what is linked to by the |event_id| and |event_type| fields.  Each
section title has a string, which is the value of the |name| field.

\section{Start/Finish Events}

Two bookend events are sent, one when BlueJ starts (before it has loaded up a
project or done anything) and one when BlueJ exits (after it has closed all
projects and finished everything).  The data will likely contain more start
events than finish events, for reasons explained in \myref{def:interruption}.

\subsection{\lstinline!"bluej_start"! and \lstinline!"bluej_finish"!}
\label{evt:bluej_start}
\label{evt:bluej_finish}

These two events have no extra accompanying event.  They have no |project_id|
and no |package_id|.  However, the |"bluej_start"| event will have zero or more
attached entries in the \tabref{extensions} that link back to this master
event -- these are extensions loaded from the system directories.

The |"bluej_start"| event will always create a new entry in the
\tabref{sessions} -- the installation details for this copy of BlueJ will be
recorded there.

\section{Project Events}

When a project is opened or closed in BlueJ, corresponding events will be
sent.  See \myref{def:projects} for details on how projects are
distinguished.  A project-open event will be followed by at least one
package-open event.

\subsection{\lstinline!"project_opening"!}
\label{evt:project_opening}

This event has no |package_id|.

This event will have zero or more entries in the \tabref{extensions} linking
back to this event, which represent extensions loaded from this project's directory.

\subsection{\lstinline!"project_closing"!}
\label{evt:project_closing}

This event will no attached data, and will have no |package_id|.

\section{Package Events}

When a package is opened or closed in BlueJ, a corresponding event is
generated.  Note that each window in BlueJ actually belongs to a package.

\subsection{\lstinline!"package_opening"!}
\label{evt:package_opening}

This event will have a |package_id|.

The event will have zero or more entries in the \tabref{source_histories} that
link back to this event.  These represent the state of the files as the project is
opened (as a |complete| snapshot, not a diff).

The event will link to an entry in the |source_hashes| table
(|event_type| will be |"SourceHash"|), which will contain a hash of
the source code in the package.  The exact details of the hash are
unimportant, but this hash can be used to quickly narrow down likely
duplicate projects (for example, spotting the starting projects from
the BlueJ book).

\label{tab:source_hashes}
\begin{table}[H]
\begin{center}
\caption[\lstinline!source_hashes! schema]{Schema for the \lstinline!source_hashes! table.  \intern}
\begin{tabular}{l@{\hspace{2cm}}l@{\hspace{1cm}}l}
Field Name & Type & Can Be Null?\\ \hline
\end{tabular}
\end{center}
\end{table}

\subsection{\lstinline!"package_closing"!}
\label{evt:package_closing}

This event will have a |package_id|, but no other associated data.


\section{Editing}

We do not capture all the different edit operations (e.g. paste, indent), we
only capture the resulting effect on the source code.

\subsection{\lstinline!"edit"!}
\label{evt:edit}

This event will have a |package_id|, pointing to the package of the class that
is edited.

This event is the central editing event in BlueJ.  These events are recorded
when multiple lines are edited at once (e.g. via paste, or deleting a
selection, or so on), or a single line has been edited and the caret
has now left that edited line.  An edit event is not generated by every single
keystroke.

There is currently no guarantee as to how edits are grouped together into a single edit
event (e.g. highlighting a selection and then pasting may appear as one event
or two events, auto-indent may appear as one edit or several edits, and many
other such examples), but an edit event will never span a compiler event, and
all edits will be flushed before a compile event.

The event will always have exactly one associated entry in the
\tabref{source_histories}, which will be a |"diff"| event.

\subsection{\lstinline!"file_add"!}
\label{evt:file_add}

This event will have a |package_id|, which indicates which package the file
was added to.

This event is fired whenever a class is added to the BlueJ project.  This may
be via the New Class menu item, or via the Import Class menu item, or via
version control.

The event will always have exactly one associated entry in the
\tabref{source_histories}, which will be a |"complete"| event.

\subsection{\lstinline!"file_delete"!}
\label{evt:file_delete}

This event will have a |package_id|, which indicates which package the file
was deleted from.

This event is fired whenever a class is removed from a BlueJ project.  This
may be via the interface, or via version control.

The event will always have exactly one associated entry in the
\tabref{source_histories}, which will be a |"file_delete"| event.

\subsection{\lstinline!"rename"!}
\label{evt:rename}

This event will have a |package_id|, which indicates which package the file
was renamed in.

This event is fired whenever a class is renamed in a BlueJ project.  A rename
occurs in BlueJ when a Java file is saved, and is parseable, and has a
different top-level class name to the name of the file.  At this point, BlueJ
silently renames the file to match the class name.  But note that until this
save occurs, many more edits may take place, so the rename of the file may
appear some time after the actual relevant source code edit takes place.

The event will always have exactly one associated entry in the
\tabref{source_histories}, which will be a |rename| event.

\section{Compilation}

\subsection{\lstinline!"compile"!}
\label{evt:compile}

This event will have a |package_id| if and only if all the compiled files were
in the same package.  However, having a |package_id| does not necessarily
indicate that the user asked for the whole package to be compiled.

This event occurs whenever a compilation of the user's classes is triggered.
(This does not include things like codepad interactions).

This master event will always have |event_type| set to |"CompileEvent"|, and
|event_id| will reference an entry in the \tabref{compile_events}.  (In turn,
that entry in the |compile_events| table, will be linked to by one or more
|compile_inputs| and zero or more |compile_outputs|, but see the
\tabref{compile_events} for more details.)

\subsection{Related Tables}

The |compile_events| table records whether the compilation was a success
(i.e. no errors -- but there could be warnings) or not.

\label{tab:compile_events}
\begin{table}[H]
\begin{center}
\caption[\lstinline!compile_events! schema]{Schema for the \lstinline!compile_events! table. 
}
\begin{tabular}{l@{\hspace{2cm}}l@{\hspace{1cm}}l}
Field Name & Type & Can Be Null?\\ \hline
\end{tabular}
\end{center}
\end{table}

The |compile_events| may then be linked to by entries in the |compile_inputs| and
|compile_outputs| table.

\label{tab:compile_inputs}
\begin{table}[H]
\begin{center}
\caption[\lstinline!compile_inputs! schema]{Schema for the \lstinline!compile_inputs! table. 
}
\begin{tabular}{l@{\hspace{2cm}}l@{\hspace{1cm}}l}
Field Name & Type & Can Be Null?\\ \hline
\end{tabular}
\end{center}
\end{table}

The |compile_inputs| table is really just a link table, linking one compile
event to many source files.

\label{tab:compile_outputs}
\begin{table}[H]
\begin{center}
\caption[\lstinline!compile_outputs! schema]{Schema for the \lstinline!compile_outputs! table. 
}
\begin{tabular}{l@{\hspace{2cm}}l@{\hspace{1cm}}l}
Field Name & Type & Can Be Null?\\ \hline
\end{tabular}
\end{center}
\end{table}

The |compile_outputs| table has an entry for each warning and each error that
occurred.  Note that these warnings and errors are reported from the standard Java
compiler from the JDK.  However, BlueJ may alter the format of some messages,
and will only display at most one error, and may ignore some warnings.  The
message recorded here is the altered format, and the |shown| flag indicates
whether or not this error or warning was actually displayed to the user.  The
|is_error| flag indicates if it was an error (otherwise it was a warning), the
|message| field shows the text that would be displayed to the user, and the
line and column fields indicate the error's location, as reported by the
compiler  (lines and columns begin at 1, not 0).

\section{Miscellaneous Events}

\subsection{\lstinline!"resetting_vm"!}
\label{evt:resetting_vm}

This event has no |package_id|, because the debug VM is per-project.

The event occurs when the VM is restarted due to an explicit user request (via
a menu).  Although it has the same effect as the |debugger_terminate| event,
it is recorded separately.

\subsection{\lstinline!"terminal_open"! and \lstinline!"terminal_close"!}
\label{evt:terminal_open}
\label{evt:terminal_close}

These events have no |package_id|, because there is only one terminal
per project.

These events occur when the terminal window is shown or closed.  It
may be shown either due to explicit triggered by the user (via a menu
item) or because output has been generated by the debug VM (which
causes the terminal to be shown if it is not already shown).

\section{Debugger Events}

Some, but not all, debugger events link to the |debugger_events|
table, which can store thread names.

\label{tab:debugger_events}
\begin{table}[H]
\begin{center}
\caption[\lstinline!debugger_events! schema]{Schema for the \lstinline!debugger_events! table. 
}
\begin{tabular}{l@{\hspace{2cm}}l@{\hspace{1cm}}l}
Field Name & Type & Can Be Null?\\ \hline
\end{tabular}
\end{center}
\end{table}

Other debugger events will link to the |breakpoints| table, which references a
source file and a line number (line numbers begin at 1, not 0).

\label{tab:breakpoints}
\begin{table}[H]
\begin{center}
\caption[\lstinline!breakpoints! schema]{Schema for the \lstinline!breakpoints! table.  \intern}
\begin{tabular}{l@{\hspace{2cm}}l@{\hspace{1cm}}l}
Field Name & Type & Can Be Null?\\ \hline
\end{tabular}
\end{center}
\end{table}

\subsection{\lstinline!"debugger_open"! and \lstinline!"debugger_close"!}
\label{evt:debugger_open}
\label{evt:debugger_close}

These events have no |package_id| (because the debugger is per-project).

These events occur when the debugger window is opened or closed.  The debugger
may open because the user explicitly opened the debugger (via the menu), or
because the code has hit a breakpoint.  No further data is attached.

\subsection{\lstinline!"debugger_terminate"!}
\label{evt:debugger_terminate}

This event has no |package_id| (because the debugger is per-project).

This event occurs when the user clicks the Terminate button in the debugger.
It kills the current debug VM (which will be subsequently restarted, but will
be idle), and thus all currently-executing user code.  No
further data is attached.

\subsection{\lstinline!"debugger_continue"!}
\label{evt:debugger_continue}

This event has no |package_id| (because the debugger is per-project).

This event occurs when the user clicks the Continue button in the debugger on
a thread which is currently not executing.
It makes the selected thread resume execution.  

The |event_type| field will always be |"DebuggerEvent"|, and the |event_id| will
reference an entry in the \tabref{debugger_events}.  The |debugger_events|
entry will have the |thread_name| filled in, but there will be no further information.

\subsection{\lstinline!"debugger_hit_breakpoint"!, \lstinline!"debugger_stepinto"!,\\ \lstinline!"debugger_stepover"! and \lstinline!"debugger_halt"!}
\label{evt:debugger_hit_breakpoint}
\label{evt:debugger_stepinto}
\label{evt:debugger_stepover}
\label{evt:debugger_halt}

These events have no |package_id| (because the debugger is per-project).

These events occur when the user code in a thread becomes suspended, for
different reasons.  The
|debugger_hit_breakpoint| event indicates that the suspension was due to
hitting a breakpoint, while |debugger_stepinto| and |debugger_stepover|
indicate that the thread has just completed a step request.  (If a step
request ends on a breakpoint, only the breakpoint is reported.)  The
|debugger_halt| is reported in the remaining case, that the user has
clicked the Halt button in the debugger (which pauses execution of the thread).

For these events, the |event_type| field will be always be |"DebuggerEvent"|,
and the |event_id| will reference an entry in the \tabref{debugger_events},
which will have the |thread_name| filled in.

There will also be one or more entries in the \tabref{stack_entries} which
will link to the |debugger_event| (not the master event) via the fields in
that table
|stack_entries.sub_event_type| and |stack_entries.sub_event_id|.

\subsection{\lstinline!"debugger_breakpoint_add"! and \lstinline!"debugger_breakpoint_remove"!}
\label{evt:debugger_breakpoint_add}
\label{evt:debugger_breakpoint_remove}

These events do have a |package_id|, which indicates the package that the
involved source file is in.

These events occur when a user adds/removes a breakpoint.  The |event_type|
field will be |"Breakpoint"| and the |event_id| will reference an entry in
the \tabref{breakpoints}, which details the location of the breakpoint.


\section{Codepad}

\subsection{\lstinline!"codepad"!}
\label{evt:codepad}

These events have a |package_id|, which indicates which package window the
codepad being used corresponds to -- more by accident than by design, each
package window has its own codepad.

These events occur each time that the user types something into the codepad
window, presses return and gets a response.  The codepad event links to an item in the |codepad_events| table (the
|event_type| will be |"CodepadEvent"|, and the |event_id| will reference an
entry in the |codepad_events| table).

\label{tab:codepad_events}
\begin{table}[H]
\begin{center}
\caption[\lstinline!codepad_events! schema]{Schema for the \lstinline!codepad_events! table. 
}
\begin{tabular}{l@{\hspace{2cm}}l@{\hspace{1cm}}l}
Field Name & Type & Can Be Null?\\ \hline
\end{tabular}
\end{center}
\end{table}

The |outcome| field indicates what happened when the code was
executed.  An |outcome| of |"success"| indicates that the code
successfully compiled, and executed without throwing an exception or without
exiting the VM early.  In this case, the fields |command| and |result|
will be filled in with the input and output respectively.


The |outcome| field being |"error"| indicates a compiler error occurred, or that the VM
exited during execution: the |command| and |error| fields will be filled in with the
input and error respectively.  The |outcome| field being |"exception"|
indicates an exception occurred: the |command| and |exception| fields
will be filled in with the input and exception respectively.


\section{Method Invocation}

Method invocations reference the |invocations| table.

\label{tab:invocations}
\begin{table}[H]
\begin{center}
\caption[\lstinline!invocations! schema]{Schema for the \lstinline!invocations! table. 
}
\begin{tabular}{l@{\hspace{2cm}}l@{\hspace{1cm}}l}
Field Name & Type & Can Be Null?\\ \hline
\end{tabular}
\end{center}
\end{table}

\subsection{\lstinline!"invoke_method"!}
\label{evt:invoke_method}

This event has a |package_id| field, which indicates which package window the
method was invoked from.  The method invoked is not guaranteed to be in this
package, as the user may have used the call-library-class function to invoke a
method on String, or any other class.

This event occurs when the user interactively (i.e. via a menu) invokes a
static method on a class or an instance method on an object on the object
bench, or uses the call-library-class menu function.  The |event_type| will
always be |"Invocation"| and the |event_id| will reference an entry in the
|invocations| table.  

In the |invocations| table, the |code| field will have the code that was generated
to execute the invocation.  The |result| field will have the result, which will be one of:

\begin{itemize}
\item |"success"| -- method was invoked successfully, the |bench_object_id|
  field will link to an entry in the \tabref{bench_objects} -- if the result was added to the
  object bench, i.e. if the method was a constructor -- and the the |type_name|
  field will have the type of the result.
\item |"compile_error"| -- an error occurred while compiling the generated code
  (e.g. because the user entered invalid parameters), the |compile_error|
  field will have the full message.
\item |"terminated"| -- the VM terminated during execution, e.g. due to a
  |System.exit| call.
\item |"exception"| -- an uncaught exception occurred during execution.  The
  |exception_class| field will have the class of the exception, the
  |exception_message| field will have the message, and there will be
  associated entries in the \tabref{stack_entries} linking to this event (the
  |invocations| entry, not the |master_events| entry).
\end{itemize}

\section{Inspector}

The inspector events reference the inspectors table.

\label{tab:inspectors}
\begin{table}[H]
\begin{center}
\caption[\lstinline!inspectors! schema]{Schema for the \lstinline!inspectors! table. 
}
\begin{tabular}{l@{\hspace{2cm}}l@{\hspace{1cm}}l}
Field Name & Type & Can Be Null?\\ \hline
\end{tabular}
\end{center}
\end{table}

\subsection{\lstinline!"inspector_show"! and \lstinline!"inspector_hide"!}
\label{evt:inspector_show}
\label{evt:inspector_hide}

This event has a |package_id| field, which indicates which package window the
inspection was invoked from.  This is transitive, for occasions where users
inspect an object in an inspector.

Both of these events link to an entry in the |inspectors| table: the
|event_type| field will be |"Inspector"|, and the |event_id| field will be a
reference to an entry in the |inspectors| table.

For the |inspector_show| event, there are two cases:

\begin{enumerate}
\item The inspection was on a class.  In this case, the |static_class| field
  will have the name of the class, and the other fields will be empty.
\item The inspection was on an object.  In this case, the |class_name| field
  indicates the type of the object, the |display_name| (whatever is displayed
  on the left-hand side of the title in the object inspector) will be present,
  and the |bench_object_id| will link to a |bench_object| if the inspected
  object is on the object bench (regardless of whether that was how this
  object was reached for inspection).
\end{enumerate}

In both cases, the |unique| will have some string that uniquely identifies an
inspector within the session (this is useful during collection, and is of no
use afterwards, because the record will be unique per inspector).

For the |inspector_hide| event, there will be a link back to the |inspectors|
entry that was hidden.

\section{Object Bench}

\subsection{\lstinline!"bench_get"!}
\label{evt:bench_get}

This event has a |package_id|, linking to which package the object
bench is associated with.

The event occurs when a user is inspecting something (an object,
class, or method result) and uses the Get button to transfer the
object to the object bench -- or when they drag a result object over
from the codepad.
  
The event links to an entry in the \tabref{bench_objects} (|event_type| will
be |"BenchObject"|), indicating
the object which was added to the bench.

If a test was being recorded at the time, the |bench_objects| entry
will have a link to the |tests| table, indicating which test was being recorded.

\subsection{\lstinline!"remove_object"!}
\label{evt:remove_object}

This event has a |package_id|, linking to which package the object
bench is associated with.

The event occurs when a user explicitly removes an object from the object bench.
  
The event links to an entry in the \tabref{bench_objects}, indicating
the object which was removed from the bench.

\section{Executing Tests}

\subsection{\lstinline!"run_test"!}
\label{evt:run_test}

This event has a |package_id|, linking to which package the test class
is in.

The event occurs when the user executes a test, either by using the
Test All command, or individually invoking a test.  Each test that is
run will get its own |run_test| event.

The event links to an entry in the |test_results| table.

\label{tab:test_results}
\begin{table}[H]
\begin{center}
\caption[\lstinline!test_results! schema]{Schema for the \lstinline!test_results! table. 
}
\begin{tabular}{l@{\hspace{2cm}}l@{\hspace{1cm}}l}
Field Name & Type & Can Be Null?\\ \hline
\end{tabular}
\end{center}
\end{table}

The |class_name| and |method_name| fields describe which test method was
executed.  The |run_time| field has the number of milliseconds that the test
took to execute.  The |result| field will have one of three values:

\begin{itemize}
\item |"success"| -- the test passed, no further information is supplied.
\item |"failure"| -- the test failed with a failed assertion.  The
  |exception_message| and |exception_trace| fields have more information.
\item |"error"| -- the test failed due to an exception (other than an assertion).  The
  |exception_message| and |exception_trace| fields have more information.
\end{itemize}

\section{Recording Tests}

\subsection{\lstinline!"start_test"!, \lstinline!"cancel_test"! and \lstinline!"end_test"!}
\label{evt:start_test}
\label{evt:cancel_test}
\label{evt:end_test}

This event has a |package_id|, linking to which package the test class
is in.

These events relate to recording a test.  The user performs various actions,
and at the end of the recording, code is generated in a new test method to
mirror the user's actions.  Each of these events links to an entry in the
|tests| table.

\label{tab:tests}
\begin{table}[H]
\begin{center}
\caption[\lstinline!tests! schema]{Schema for the \lstinline!tests! table. 
}
\begin{tabular}{l@{\hspace{2cm}}l@{\hspace{1cm}}l}
Field Name & Type & Can Be Null?\\ \hline
\end{tabular}
\end{center}
\end{table}

Each entry in the |tests| table is one attempt (either ongoing, successful, or
aborted) at recording a test.  The |source_file_id| links to the source file
in which the final test method (named according to |method_name|) will be
generated.  The |status| field will be |"Recording"|, |"Recorded"| or
|"Cancelled"|, according to the test's current status.

The |session_subid| field is an internal identifier used by the client to
identify tests, and can be ignored.  The |session_id| field links to the
session in which the test was recorded.

\subsection{\lstinline!"assertion"!}
\label{evt:assertion}

This event has a |package_id|, linking to which package the test was
recorded for.

The event links to an entry in the |assertions| table.

\label{tab:assertions}
\begin{table}[H]
\begin{center}
\caption[\lstinline!assertions! schema]{Schema for the \lstinline!assertions! table. 
}
\begin{tabular}{l@{\hspace{2cm}}l@{\hspace{1cm}}l}
Field Name & Type & Can Be Null?\\ \hline
\end{tabular}
\end{center}
\end{table}

The |assertions| table links to an entry in the |invocations| table that the assertion
relates to, and to an entry in the \tabref{tests} for the test which
this assertion will be recorded into.

The |assertion| field is the JUnit function that will be recorded,
e.g. |"assertEquals"|.  The |param1| field and |param2| fields are
the other parameters to the assertion.

\subsection{\lstinline!"bench_to_fixture"! and \lstinline!"fixture_to_bench"!}
\label{evt:bench_to_fixture}
\label{evt:fixture_to_bench}

These events have a |package_id|, linking to the package in which the object bench
and test class are.

These events occur when the user captures all the objects on the object bench
and turns their creation into a test setup phase (bench to fixture) or runs a
test setup to put the objects back onto the object bench (fixture to bench).

These events link to an entry in the |fixtures| table.

\label{tab:fixtures}
\begin{table}[H]
\begin{center}
\caption[\lstinline!fixtures! schema]{Schema for the \lstinline!fixtures! table. 
}
\begin{tabular}{l@{\hspace{2cm}}l@{\hspace{1cm}}l}
Field Name & Type & Can Be Null?\\ \hline
\end{tabular}
\end{center}
\end{table}

Each entry in the |fixtures| table links to the source file holding the test
class involved.  A new fixture will be created for each |"bench_to_fixture"|
or |"fixture_to_bench"| event.

\label{tab:bench_objects_fixtures}
\begin{table}[H]
\begin{center}
\caption[\lstinline!bench_objects_fixtures! schema]{Schema for the \lstinline!bench_objects_fixtures! table. 
}
\begin{tabular}{l@{\hspace{2cm}}l@{\hspace{1cm}}l}
Field Name & Type & Can Be Null?\\ \hline
\end{tabular}
\end{center}
\end{table}

The |bench_objects_fixtures| table is a simple link table that links object
bench objects to test fixtures.  In the case of |"bench_to_fixture"|, these
objects are already on the bench and will be removed by this event.  In the
case of |"fixture_to_bench"|, these objects will be added to the bench, and
are created by this event.

\section{Version Control}

\label{tab:version_control_events}
\begin{table}[H]
\begin{center}
\caption[\lstinline!version_control_events! schema]{Schema for the \lstinline!version_control_events! table. 
}
\begin{tabular}{l@{\hspace{2cm}}l@{\hspace{1cm}}l}
Field Name & Type & Can Be Null?\\ \hline
\end{tabular}
\end{center}
\end{table}

\label{tab:version_control_files}
\begin{table}[H]
\begin{center}
\caption[\lstinline!version_control_files! schema]{Schema for the \lstinline!version_control_files! table. 
}
\begin{tabular}{l@{\hspace{2cm}}l@{\hspace{1cm}}l}
Field Name & Type & Can Be Null?\\ \hline
\end{tabular}
\end{center}
\end{table}

Each version control event has an associated entry in the
|version_control_events| table, with a |vcs_type| (built-in options:
|"CVS"|, |"SVN"|) and a |protocol| (e.g. |"http"| or |"svn+ssh"|).

\subsection{\lstinline!"vcs_history"! and \lstinline!"vcs_share"!}
\label{evt:vcs_history}
\label{evt:vcs_share}

These events get an associated entry in the |version_control_events|
table, but no more data than that.

\subsection{\lstinline!"vcs_commit"! and \lstinline!"vcs_update"!}
\label{evt:vcs_commit}
\label{evt:vcs_update}

These events get an associated entry in the |version_control_events|
table, and also zero or more entries in the |version_control_files|
table, with just the |file| field filled in.

\subsection{\lstinline!"vcs_status"!}
\label{evt:vcs_status}

This event gets an associated entry in the |version_control_events|
table, and also zero or more entries in the |version_control_files|
table, with the |file| field filled in, and the |status| field filled
in.

\chapter{Other Database Tables}

\section{Users}

\label{tab:users}
\begin{table}[H]
\begin{center}
\caption[\lstinline!users! schema]{Schema for the \lstinline!users! table. 
}
\begin{tabular}{l@{\hspace{2cm}}l@{\hspace{1cm}}l}
Field Name & Type & Can Be Null?\\ \hline
\end{tabular}
\end{center}
\end{table}

The |uuid| field is the unique user identifier, as discussed in \myref{def:users}.
The |created_at| field shows the time that this user record was
created (the first time their UUID was seen), and the |last_update| field shows the last time that data was entered for
this user.

\section{Sessions}

\label{tab:sessions}
\begin{table}[H]
\begin{center}
\caption[\lstinline!sessions! schema]{Schema for the \lstinline!sessions! table. 
}
\begin{tabular}{l@{\hspace{2cm}}l@{\hspace{1cm}}l}
Field Name & Type & Can Be Null?\\ \hline
\end{tabular}
\end{center}
\end{table}

Each session is associated with a specific user.  There is a link to
the \tabref{installation_details} -- specifically, the entry that
accompanied the |bluej_start| event that began this session.  The
|identifier| and |last_sequence_id| fields are used for internal
book-keeping, and can be ignored.  There is a link to an entry in the
\tabref{participants} if this user had participant information filled in
during the last event that was added to this session.

\section{Installation Details}

\label{tab:installation_details}
\begin{table}[H]
\begin{center}
\caption[\lstinline!installation_details! schema]{Schema for the \lstinline!installation_details! table.  \intern}
\begin{tabular}{l@{\hspace{2cm}}l@{\hspace{1cm}}l}
Field Name & Type & Can Be Null?\\ \hline
\end{tabular}
\end{center}
\end{table}

The |installation_details| table contains the BlueJ |interface_language| -- the
possible options are listed in the bluej.defs file.  The |operating_system|
field glues together (with slashes) the output of Java's |System.getProperty|
call for the keys |"os.name"|, |"os.arch"| and |"os.version"|.  The
|bluej_version| is the version of BlueJ they are running, and |java_version|
is the Java version being used to run BlueJ.

\section{Extensions}

\label{tab:extensions}
\begin{table}[H]
\begin{center}
\caption[\lstinline!extensions! schema]{Schema for the \lstinline!extensions! table. 
}
\begin{tabular}{l@{\hspace{2cm}}l@{\hspace{1cm}}l}
Field Name & Type & Can Be Null?\\ \hline
\end{tabular}
\end{center}
\end{table}

The |extensions| table contains the (self-reported) name and version of a
BlueJ extension.  BlueJ comes with the submitter extension installed by default.

\section{Projects}

\label{tab:projects}
\begin{table}[H]
\begin{center}
\caption[\lstinline!projects! schema]{Schema for the \lstinline!projects! table. 
}
\begin{tabular}{l@{\hspace{2cm}}l@{\hspace{1cm}}l}
Field Name & Type & Can Be Null?\\ \hline
\end{tabular}
\end{center}
\end{table}

Each project is associated with a specific user.  A project is stored with the
md5 hash of its full path (|path_hash|), to try to identify the project across multiple
sessions.  For anonymity purposes, the full path is not sent to the server
(because it will typically have the user-name in the path), but the directory
name of the project (|name|) is sent, as this may be informative.

\section{Packages}

\label{tab:packages}
\begin{table}[H]
\begin{center}
\caption[\lstinline!packages! schema]{Schema for the \lstinline!packages! table. 
}
\begin{tabular}{l@{\hspace{2cm}}l@{\hspace{1cm}}l}
Field Name & Type & Can Be Null?\\ \hline
\end{tabular}
\end{center}
\end{table}

Each package belongs to a specific project.  A package is identified by its
fully-qualified package name.  The default, unnamed package will have the
empty string as its name.

\section{Source Files}

\label{tab:source_files}
\begin{table}[H]
\begin{center}
\caption[\lstinline!source_files! schema]{Schema for the \lstinline!source_files! table. 
}
\begin{tabular}{l@{\hspace{2cm}}l@{\hspace{1cm}}l}
Field Name & Type & Can Be Null?\\ \hline
\end{tabular}
\end{center}
\end{table}

Each entry in the |source_files| table corresponds to a particular
source path in a project.  (That is, if you delete a source file and
generate a new file with the same name, it may share the entry of the
original file.)

Each source file will have its path relative to the base of the
project.  So files that are within packages will have a path with the
parent package names as directories (which may use forward slash or
backslash for separation).

\section{Source Histories}

\label{tab:source_histories}
\begin{table}[H]
\begin{center}
\caption[\lstinline!source_histories! schema]{Schema for the \lstinline!source_histories! table. 
}
\begin{tabular}{l@{\hspace{2cm}}l@{\hspace{1cm}}l}
Field Name & Type & Can Be Null?\\ \hline
\end{tabular}
\end{center}
\end{table}

Each entry in the |source_histories| table belongs to a particular
|master_event|, and is associated with a given |source_file|.  Based on the
|source_history_type| field, there can be different meanings to the |content| field:

\begin{itemize}
\item |source_history_type| is |"complete"| -- |content| has the complete text
  of the source file (anonymised).  This will be
  used by file additions, and source histories added on loading a package.
\item |source_history_type| is |"diff"| -- |content| has the diff (in unified diff format) from
  the previous version of the file.
\item |source_history_type| is |"rename"| -- |content| holds the name of the file before the rename.  During the
  rename, the source history will be updated to feature the new name of the
  file.
\item |source_history_type| is |"file_delete"| -- |content| will be blank.
\end{itemize}

\section{Stack Entries}

\label{tab:stack_entries}
\begin{table}[H]
\begin{center}
\caption[\lstinline!stack_entries! schema]{Schema for the \lstinline!stack_entries! table. 
}
\begin{tabular}{l@{\hspace{2cm}}l@{\hspace{1cm}}l}
Field Name & Type & Can Be Null?\\ \hline
\end{tabular}
\end{center}
\end{table}

Any time that an event is associated with a stack trace, it will have
zero or more entries in this table: one for each line in the stack
trace.  (Some stack traces may be filtered by BlueJ -- the recorded
trace will always be the filtered version, i.e. the version that the
user actually sees.)

The |entry| field will be the index within the stack trace.  The top
line (typically, the actual location of the thrown exception, or
breakpoint, etc) of the trace will have |entry| of zero, with an
increasing |entry| index for each subsequent line.

The |class_source_name| field will be the name of the file involved,
and |line_number| will be the line number (line numbers start at 1).
So if an exception is thrown in an anonymous inner class of class Foo,
this field will be |"Foo.java"| -- even if Foo is inside a package.
The source name loses all package information, but this can be
recovered via the |class_name| field, which is the fully qualified
name of the class at the specific source location.

The |sub_event_id| and |sub_event_type| fields are what is used to
link the stack entries to their parent event -- see documentation
throughout the manual for events that have associated stack traces.
\section{Bench Objects}

\label{tab:bench_objects}
\begin{table}[H]
\begin{center}
\caption[\lstinline!bench_objects! schema]{Schema for the \lstinline!bench_objects! table. 
}
\begin{tabular}{l@{\hspace{2cm}}l@{\hspace{1cm}}l}
Field Name & Type & Can Be Null?\\ \hline
\end{tabular}
\end{center}
\end{table}

Objects on the BlueJ object bench get an entry in the |bench_objects|
table.  The |name| field is the name/label of the object, and
|class_name| is the type of the object.  The |package_id| field links
to the package which is associated with the object bench that this
object is on.  The |session_id| field links to the session during
which this object existed.

The |test_id| field links to an entry in the \tabref{tests} -- this
link is present if (and only if) the given test was being recorded
when this object was added to the bench.

\section{Participant Information}
\label{sec:views}

To keep participant information (see \myref{sec:participants}) hidden,
we hide the actual table involved, |participants|, and instead offer three
database views of the information.

\label{tab:participant_identifiers_for_experiment}
\begin{table}[H]
\begin{center}
\caption[\lstinline!participant_identifiers_for_experiment! schema]{Schema for the \lstinline!participant_identifiers_for_experiment! table. 
}
\begin{tabular}{l@{\hspace{2cm}}l@{\hspace{1cm}}l}
Field Name & Type & Can Be Null?\\ \hline
\end{tabular}
\end{center}
\end{table}

The |participant_identifiers_for_experiment| view allows you to see the
available participant identifiers for a given experiment identifier.  You can
call the view using a statement such as:

\begin{lstlisting}
select p.participant_identifier from
  (select @experiment:="myexperiment") unused,
  participant_identifiers_for_experiment p;
\end{lstlisting}

The inner select is a trick to set the |@experiment| parameter to be the
experiment identifier.  The view then retrieves a list of participant
identifiers based on the experiment identifier that you supply.


\label{tab:sessions_for_experiment}
\begin{table}[H]
\begin{center}
\caption[\lstinline!sessions_for_experiment! schema]{Schema for the \lstinline!sessions_for_experiment! table. 
}
\begin{tabular}{l@{\hspace{2cm}}l@{\hspace{1cm}}l}
Field Name & Type & Can Be Null?\\ \hline
\end{tabular}
\end{center}
\end{table}

The |sessions_for_experiment| view allows you to retrieve information for all
the sessions associated with a given experiment identifier.  A session is
associated with an experiment identifier (and similarly with a participant
identifier) if during the last event logged in that session, the user had
their experiment identifier set to the given value.  So if a user sets their
identifiers partway through a session (and does not change it again before the
end of the session), the whole session is deemed to be part of the experiment.

The columns returned are those from the \tabref{sessions}, plus the
|participant_identifier| column.

The view can be called using a statement such as:

\begin{lstlisting}
select s.* from
  (select @experiment:="myexperiment") unused,
  sessions_for_experiment s;
\end{lstlisting}

\label{tab:sessions_for_experiment_participant}

% The |sessions_for_experiment_participant| view is very similar to
% |sessions_for_experiment|, but also allows you to specify the
% |participant_identifier|.  The columns returned are those from the
% \tabref{sessions}, and the view can be called using:

% \begin{lstlisting}
% select s.* from
%   (select @experiment:="myexperiment", @participant:="p001") unused,
%   sessions_for_experiment_participant s;
% \end{lstlisting}

\chapter{Data Access}

The data collection project involves two machines, running Ubuntu.  Each
machine has two 6-core hyperthreaded 2.5Ghz Xeon processors, with 32GB of RAM,
and four 2TB drives set up in RAID 5 configuration to hold the database.

One machine, nicknamed black, is
responsible for collecting the data.  It solely runs the MySQL database, and
the Ruby on Rails server.  

The second machine, nicknamed white, has a live-cloned database (named |blackbox_production|) from the recording
machine.  White is the machine to which you will have an account, and be able
to log in and perform data analysis.  The intention is that this way, nothing
that anyone does while examining the database on white can affect the
recording on black.

\section{SSH Access}

The white machine can be accessed by SSHing to |white.kent.ac.uk| on the
standard port (22).  The machine runs Ubuntu.  If you need any software
installed, then we can install any standard Ubuntu packages for you (email
|blackbox-admin@bluej.org|).

\section{MySQL}

The white machine runs MySQL on the standard port (3306) on the localhost
interface.  If you want to access the database directly from your own machine,
you can use an SSH tunnel:

\begin{lstlisting}
ssh -L 3306:localhost:3306 ssh_username@white.kent.ac.uk
\end{lstlisting}

While that SSH session is open, you can access the database by running
|mysql -u username -p blackbox_production|.  You will have been given
the MySQL details separately when you got an account on the machine.

\section{Feed}

To get familiarity with the data collection, we provide a mechanism to observe
a live feed of your own data.  The live feed is available, from the
white machine, on |http://localhost:3080/feed|.  It is probably easiest to
port-forward 3080 to your local machine using:

\begin{lstlisting}
ssh -L 3080:localhost:3080 ssh_username@white.kent.ac.uk
\end{lstlisting}

Then access |http://localhost:3080/feed| on your own machine.  You can find
your own UUID in your |bluej.properties| file -- look for the line
|blackbox.uuid=...|.  If you paste your UUID into the feed page and click the
button, it will refresh every five seconds until closed (be warned!) to show
you the data for that UUID.  Looking at this webpage while running BlueJ
should give a good feel for what is being recorded.

\section{Analysis Process}

We anticipate that researchers will do some of their analysis on the white
machine, and some on their own machine.  The advantage of using the white
machine is that it is high-powered, and has immediate access to all of the
data, without needing to copy the data across the Internet.  The disadvantage
is that it can be a bit restrictive to operate on a remote machine for all
analysis.  Researchers may therefore want to do a split, with some
pre-processing on the white machine, followed by more investigation on their
own machine.

\section{Data Snapshot}

When researchers publish their data, they will specify their data source.  It
is worth remembering that the Blackbox data is continually being added to, so
it is also necessary to specify which time period data was used from.  We
suggest choosing a cut-off date (e.g. the beginning date of the analysis) and
ensuring that only data up until that cut-off date is used, so that the source
of the data can be correctly specified as ``the Blackbox data repository was
used, up until 2013-03-31''.  You can use the date stamps in the
\tabref{master_events} to decide which data to use in your study.  (Bear in
mind that this may bisect some sessions -- make sure you do not process the
events in a session which fall outside your cut-off date).

\chapter{Data Use Policies}

The data that we collect is not public.  This is part of our ethics approval
-- it is only visible to us, and selected researchers (i.e. those to whom we
give an account on the white machine).  

\section{Data Caching/Copying}

We realise that researchers will ultimately end up taking some data (probably
processed) on to their own machine, as a pragmatic step during the analysis.
We ask that you act responsibly, and make sure that the raw data is not made
public, and also that you take care to not let your account on the white
machine be compromised.  If you need multiple researchers at your institution
to have access, we would prefer giving them all accounts rather than
encouraging the practice of sharing accounts (and thus passwords).


\section{Data Use}

The data is intended for education research, and is anonymous.  We expect you
to not attempt to identify any students or data within the database, except
your own, and that of any participants in a study for which you have ethical
approval.  In particular, we forbid any use of the data for the purposes of
plagiarism detection (and similar) in students at your local institution.  The
participation of the users in the experiment is conditional on the data being
anonymous, and should not have any negative impacts for the participants.

\section{Sharing and Publication}

Researchers should not share the data with anyone -- any other researchers
wanting access should ask us for accounts.  Obviously, at some point,
researchers may want to publish their findings.  We expect that the data will
either be published in an aggregated form (which thus retains anonymity), and
that any individual case studies/vignettes will be manually anonymised, if
necessary.

\chapter{Beta}

\section{Aims}

The aim of the beta is two-fold:

\begin{itemize}
\item Firstly, we want to test the software outside of our local setup.  We
  would appreciate if you were able to test the software in a variety of
  settings, e.g. from your home machine, from a student machine in your labs.
  This will help pick up any technical problems with slow connections, or
  firewalls, or unexpected lab setups and so on.
\item Secondly, we invite feedback on the data schema (described in
  this document).  Is there something we're not recording that you want to
  know?  Are we recording something in a way that will inhibit easy
  analysis?  Are some of our caveats too big, or our guarantees too
  weak?  The recording mechanism is much easier to alter before the
  public release than afterwards.
\end{itemize}

We also welcome feedback on this document: is there anything that is
not clear, or missing, or that could be improved?

Send comments either to the mailing list |blackbox@bluej.org| or
directly to us at |blackbox-admin@bluej.org|.

\section{Artificialities}

\begin{itemize}
\item UUID is shown in preferences dialog
\item Message logged to debug log if there is a problem with event submission
\end{itemize}

\section{Notes}

During the beta, it may prove necessary to change the data schema.  In this
case, data may be wiped, and analysis tools may need adjusting.  This is in
order to provide the best setting for collecting the actual data: once the
beta is over, there will be a stable interface to code against, and no data
will be removed.

\end{document}
